\section{Open Questions}

The replication surfaced five substantive open questions that would extend or stress-test the
Fukui et al.\ IMK + SLDR + ALDH$^{+}$ framework. Each is stated with the observational basis from
the replication, and concrete next steps runnable on free / no-outreach compute.

\subsection*{Q1. Does the ALDH$^{+}$ $\leftrightarrow$ $f_s$ correspondence generalize to other stem-cell markers?}

\textbf{Basis.} The IMK model treats $f_s$ as a scalar stem-cell fraction anchored to ALDH$^{+}$
aldefluor measurements. Claim~A (4/4 agreement within $1\sigma$) is impressive but ALDH activity is
a metabolic assay, not a lineage marker, and overlaps only partially with CD44$^{+}$/CD133$^{+}$
in oral SCC. The IMK derivation nowhere requires the marker be ALDH.

\textbf{Next steps.} Refit IMK against published CD44$^{+}$/CD133$^{+}$/side-population sort +
clonogenic data for the same or comparable oral SCC lines. Compare $f_s$ posteriors
marker-by-marker. Marker-invariance would strengthen a true biological-subpopulation
interpretation; marker-dependence would reduce $f_s$ to a phenomenological knob.

\subsection*{Q2. Can the lumped $(a+c)$ SLDR rate be replaced by an explicit repair-pathway model
(NHEJ + HR + MMEJ)?}

\textbf{Basis.} $(a+c)$ is a single lumped repair-rate parameter ($\sim 1.3$--$2.8$~h$^{-1}$).
Modern DDR models (MEDRAS, LEM-IV, McMahon--Prise 2016) decompose DSB repair into distinct
pathways with cell-cycle-dependent kinetics. It is unclear whether the $\sim 2\times$
$w_{\mathrm{SLDR}}$ increase in HSC2-R reflects a genuinely upregulated pathway or is
degeneracy in a lumped fit.

\textbf{Next steps.} Wire SLDR into a two-pathway repair model (fast NHEJ + slow HR) with rates
from Reynolds et al.\ 2012 or MEDRAS defaults. Refit Fig.~5 + Fig.~6. Compare AIC/BIC vs.\ the
lumped $(a+c)$. If the two-pathway model wins, decompose the $w_{\mathrm{SLDR}}$ increase into
NHEJ-fold-change vs.\ HR-fold-change; cross-reference RAD51/53BP1 foci data.

\subsection*{Q3. How sensitive is the fit to the operational definition of ALDH$^{+}$?}

\textbf{Basis.} Fig.~3 reports ALDH$^{+}$ from $0.97 \pm 0.68\%$ (SAS) to $12.61 \pm 6.11\%$
(HSC2-R). Aldefluor-gate placement (top 0.5\% vs.\ top 1\% vs.\ top 2\% of DEAB-inhibited signal)
is a defensible choice with material effect. Reported $\sim \pm 50\%$ error bars suggest
substantial gating sensitivity that the current agreement does not disentangle.

\textbf{Next steps.} Reprocess flow data (if released, or via comparable FCS files on
FlowRepository) under three gating protocols: (a) fixed 1\% DEAB threshold, (b)
2$\sigma$-above-DEAB-mean, (c) automated GMM. Rerun IMK MCMC with each ALDH$^{+}$ prior. Report
$f_s$ posterior shift per gate. Quote gating-uncertainty envelope explicitly.

\subsection*{Q4. Do the model's predictions transfer to other tumor types
(glioblastoma, breast, colon)?}

\textbf{Basis.} The paper is scoped to oral SCC. ALDH$^{+}$ radioresistance is also reported in
GBM (Bao), breast (Ginestier), and colon (Huang), with different baseline ALDH$^{+}$ fractions.
It is unknown whether IMK + $f_s$ is universal or oral-SCC-specific; $w_{\mathrm{SLDR}} \sim 1.9$
for HSC2-R is not obviously predicted by any general theory.

\textbf{Next steps.} Assemble 5--8 published cross-tumor cell-line pairs (parental + resistant
sub-line) covering GBM (U87, U251), breast (MCF-7, MDA-MB-231), pancreatic (Panc-1). Extract
survival + ALDH$^{+}$ from published figures. Refit IMK globally with shared baseline and
per-tumor-type offset. Test whether $w_{\mathrm{SLDR}}$ clusters by tissue or by baseline
ALDH$^{+}$ fraction.

\subsection*{Q5. Can radiosensitivity be predicted from ALDH-isoform transcriptomic expression
alone, obviating wet-lab clonogenic assays?}

\textbf{Basis.} IMK fitting currently requires days of clonogenic assay per cell line. If ALDH
isoform expression (ALDH1A1, ALDH1A3, ALDH3A1) from single RNA-seq correlated with fit
parameters, the model could be used prospectively on patient biopsies. Depmap has RNA-seq +
radiation response for many of these lines.

\textbf{Next steps.} Pull Depmap CCLE RNA-seq for oral-SCC lines (SAS, HSC-2, CAL-27, HSC-3,
HSC-4, Ca9-22 etc.) plus reported clonogenic $\alpha/\beta$. Train leave-one-out ridge
regression: $\log(w_{\mathrm{SLDR}}) \sim$ ALDH1A1 $+$ ALDH1A3 $+$ ALDH3A1 $+$ BRCA1 $+$ RAD51
$+$ XRCC5. Report cross-validated MSE and Shapley values per gene. If ALDH isoforms dominate,
this validates the mechanistic claim; if repair genes dominate, $w_{\mathrm{SLDR}}$ is really a
repair-machinery readout. Extend to TCGA HNSC as a second validation cohort.
